% Nouveau EGA V (cote 25): archive cote 25, site page 99 (PDF page 100).
% typescript. Draft compared: none.
% AI-transcribed by Claude Opus 5.5 (claude-opus-5-5, adaptive thinking, effort high), Message Batches API, batch ega25b (msgbatch_01A7n5YuXubUPNK1zsepi8RG), 2026-11-13; checked against both supplied scan parts; see CHECK_DECISIONS.tsv.
\scanpage{25/99}
% NOTE: two long diagonal pen strokes run across the whole page: one from above "deuxième" (line 1) down to the left to "Corollaire" (line 5), and one from there down to the foot of the page. The whole page appears struck through by hand.
\mspage{2}

déduit de $f:P\hand{\longrightarrow} S$, le deuxième  déterminé par l'élément  cl $(\underline{O}_P(1))$ de
Pic(P). On en conclut un homomorphisme canonique
(x) \qquad Pic(S)x $\underline{Z} \hand{\longrightarrow}$ Pic(P)  ,
défini d'ailleurs sans hypothèse restrictive sur S ni sur $\underline{E}$.  Ceci posé:
\uline{Corollaire} 1.4.\placedmargin{\ensuremath{\swarrow}}   Sous les conditions de 1.\unsure{1|4}. , si $S\neq\emptyset$ 
 alors l'homomorphisme précédent est  injectif , et même
bijectif lorsque S est connexe\hand{.}

\placedmargin{[} Lorsqu'on  abandonne l'hypothèse
sur le rang $\hand{\geqslant} 2$,  il résulte de 1.2. et 1.\unsure{4|3}. que l'homomorphisme précédent
si S est connexe, mais
est encore surjectif  pas nécessairement \unsure{injectif|bijectif}, le noyau étant
isomorphe 
à $\underline{Z}$  resp. à  Pic(S)x$\underline{Z}$  lorsque  $\underline{E}$  est de rang 1\hand{,} resp. de rang 0.

Démontrons  1.1.  en commençant par l'unicité.  Tout d'abord, si S est
le spectre d'un corps, notons que  $\underline{O}(n)$ n'est isomorphe à $\underline{O}(m)$ i.e.
$\underline{O}(m-n)$  n'est  isomorphe à $\underline{O}_P$ , que si  n=m i.e.  n-m=0 .  Cela provient
du fait que $\underline{O}(1)$ est ample et  dim P $\hand{\geqslant}$1 (d'après l'hypothèse rang E $\hand{\geqslant}$2):
on peut en effet supposer  d = m-n $\hand{\geqslant}$0, si on avait d $\hand{>}$0, alors $\underline{O}(d)$
serait ample, et ne pourrait être isomorphe à  $\underline{O}_P$ que si P est quasi-affi-%
ne, donc fini (puisqu'il est propre sur k).  Ceci prouve déjà  l'unicité
de la famille $(S_n)_{n\hand{\in}\underline{Z}}$  considérée dans 1.1.  Pour l'unicité de $\underline{M}$ à isom.
unique près, on est ramené au cas où  $S=S_n$ , je dis que dans ce cas on a
un isomorphisme  (uniquement déterminé en termes de l'isom  $\underline{L} \hand{\xrightarrow{\sim}} f$$^*$$(\underline{M})(n)$):
\[ \text{(xx)} \qquad M \hand{\xrightarrow{\sim}} f_*(\underline{\unsure{L|E|M}}(-n)) \quad . \]
En effet,  l'isomorphisme  $\underline{L} \hand{\xrightarrow{\sim}} f^*(\underline{M})(n)$   définit un isomorphisme
$\underline{L}(-n) \hand{\xrightarrow{\sim}} f^*(\underline{M})$  et par suite un isomorphisme du deuxième membre de (xx)
avec $f_*(f^*(\underline{M}))$, qui lui-même est isomorphe \unsure{(|à}\,$\underline{M}$ étant loc libre) à
$\underline{M}\otimes f_*(\underline{O}_P)$ , or  (réf)  $f_*(\underline{O}_P) \hand{\xleftarrow{\sim}} \underline{O}_{\unsure{S|S_S}}$ ,  d'où l'\unsure{isomorphisme|isemorphisme} (xx).\placedmargin{\ensuremath{\backslash}}

\placedmargin{93}
